Contoh: Puji syukur ke hadirat Allah SWT atas limpahan rahmat, karunia, serta petunjuk-Nya sehingga tugas akhir berupa penyusunan skripsi ini telah terselesaikan dengan baik. Dalam hal penyusunan tugas akhir ini penulis telah banyak mendapatkan arahan, bantuan, serta dukungan dari berbagai pihak. Oleh karena itu pada kesempatan ini penulis mengucapkan terima kasih kepada:


\begin{enumerate}
	\item{Bapak Prof. Ir. Hanung Adi Nugroho, S.T., M.Eng., Ph.D., IPM., SMIEEE., selaku Ketua Jurusan Teknik Elektro dan Teknologi Informasi Fakultas Teknik Universitas Gadjah Mada,}
	
	\item{Bapak Dr.Eng. Ir. Igi Ardiyanto, S.T., M.Eng., IPM., ASEAN Eng., SMIEEE., selaku dosen pembimbing pertama yang telah memberikan banyak bantuan, bimbingan, serta arahan dalam Tugas Akhir ini,}
	\item{Bapak Ir. Agus Bejo, S.T., M.Eng., D.Eng., IPM., selaku dosen pembimbing kedua yang juga telah memberikan banyak bantuan, bimbingan, serta arahan dalam Tugas Akhir dan kegiatan-kegiatan yang lain,}
	\item{Bapak Widyawan, S.T., M.Sc., Ph.D., selaku dosen pembimbing akademik yang telah memberikan pendampingan dan arahan selama masa perkuliahan,}
	\item{Seluruh Dosen di Jurusan Teknik Elektro dan Teknologi Informasi FT UGM, yang tidak bisa disebutkan satu-satu, atas ilmu dan bimbingannya selama penulis berkuliah di JTETI,}
	\item{Ibu dan Bapak yang selama ini telah sabar membimbing, mengarahkan, dan mendoakan penulis tanpa kenal lelah untuk selama-lamanya}

\end{enumerate}

Akhir kata penulis berharap semoga skripsi ini dapat memberikan manfaat dalam bidang teknologi informasi untuk penelitian selanjutnya di masa mendatang.

\begin{flushright}
	\begin{tabular}{c}
		Yogyakarta, 25 Juni 2026 \\
		\vspace{1cm} \\
		Penulis
	\end{tabular}
\end{flushright}